\documentclass[10pt]{article}
\usepackage[margin=0.8in]{geometry}
\usepackage[T1]{fontenc}
\usepackage{lmodern}
\usepackage{amsmath,amssymb}
\usepackage{booktabs}
\usepackage{enumitem}
\usepackage{titlesec}
\usepackage[hidelinks]{hyperref}
\usepackage{microtype}

\setlist{nosep,leftmargin=*}
\titlespacing*{\section}{0pt}{6pt}{3pt}
\titlespacing*{\subsection}{0pt}{4pt}{2pt}
\titleformat*{\section}{\large\bfseries}
\titleformat*{\subsection}{\normalsize\bfseries}
\setlength{\parskip}{3pt}
\setlength{\parindent}{0pt}

\title{\vspace{-2.2em}\textbf{In-Context Engram: An Attention-Free Exact Copy Path\\for Hybrid Language Models with Less KV Cache}}
\author{Phung Gia Huy \quad \texttt{phunghuysoict@gmail.com}}
\date{Project Proposal --- October 2026}

\begin{document}
\maketitle
\vspace{-2.5em}

\section{Problem Description \& Current Limitations}

Production LLMs increasingly use \emph{hybrid} architectures: about three linear-attention layers (e.g.\ Gated DeltaNet, GDN) for every full-attention layer, as in Qwen3.5, Qwen3-Next and Kimi Linear. Linear layers keep a fixed-size recurrent state, so only the full-attention layers keep a KV cache that grows with context. At long context that KV cache dominates memory and decode time. The obvious fix, fewer full-attention layers, breaks one ability first: \textbf{exact recall of earlier text} (identifiers, numbers, code, quotes). Recurrent models cannot copy strings longer than their state can hold~[1], and in hybrids the long-range retrieval is done almost entirely by the full-attention layers~[2,3].

Much of that retrieval is \emph{exact copying}. In Zoology~[4], 82\% of the quality gap between attention and attention-free models on the Pile is explained by in-context recall. Fewer than 5\% of heads are ``retrieval heads'' that copy-paste from context, and removing them breaks retrieval~[5]; only these heads need the full KV cache~[6]. Yet exact copying needs no learned vectors: if the last $n$ tokens appeared before, the token that followed them is a strong guess for the next token, and a hash table finds it in $O(1)$ from token IDs alone.

\textbf{Limitation.} Current models pay for exact copying with the most expensive resource they have, a per-token KV cache over all past tokens. Existing ways to cut it (KV compression, head-level hybrids~[7]) still store vectors and fail first on verbatim recall. \textbf{Our thesis:} attention is not bad at copying, it is an \emph{expensive} way to copy.

\section{Technical Proposal}

\subsection{Approach: In-Context Engram (ICE)}

\textbf{Exact lookup (non-learned).} At each position $t$ and for each order $n\in\{1,2,3,4,6,8\}$, find the most recent $p<t$ with $x_{p-n+1..p}=x_{t-n+1..t}$ and return the \emph{successor} $s_n(t)=x_{p+1}$, the match length and the distance $t-p$. Since $p+1\le t$, ICE never sees the future. It stores only token IDs and one hash-table entry per order (${\approx}100$ bytes/token, versus ${\approx}2$--$3$\,KB/token for a single full-attention layer), and costs $O(1)$ per token. In training, features are computed in the data loader on CPU with a vectorized sort-based routine.

\textbf{Learned use.} ICE gives \emph{hints}, not rules; the model learns how much to trust them:
\begin{itemize}
  \item \emph{Input side (after layer 2):} $h_t \leftarrow h_t + \sum_n g_n(h_t)\, f_n$, with $f_n = W_n E[s_n(t)] + L[\text{len}] + D[\text{dist}]$ and a context-dependent gate $g_n=\sigma\!\big(\langle \mathrm{RMSNorm}(h_t),\mathrm{RMSNorm}(f_n)\rangle/\sqrt d\big)$.
  \item \emph{Output side:} a learned copy bonus $\mathrm{logits}_t[s^*(t)] \mathrel{+}= \mathrm{softplus}(w^\top h_t + b_{\text{len}})$ on the longest-match successor $s^*(t)$.
  \item \emph{ICE-dropout:} ICE is disabled for 20\% of sequences, so attention still learns fuzzy retrieval (guarding against ``lookup laziness'').
\end{itemize}

\textbf{Hypotheses.} (H1) ICE closes most of the copy-task gap between a 7:1 and a 3:1 hybrid. (H2) 7:1\,+\,ICE matches 3:1 on LM loss. (H3) ICE \emph{adds} to head-level hybrids~[7] at equal KV budget, i.e.\ it is complementary rather than an alternative. (H4, risk) Without ICE-dropout, paraphrased retrieval degrades.

\subsection{Evaluation Methods}

\textbf{Step 0 --- motivation check (no training).} On Qwen3.5-0.8B, compute per-token loss with full attention and with its full-attention layers restricted to a 512-token window. Split the loss increase by whether the ICE successor is correct. If most of the increase falls on ICE-correct tokens, long-range attention is mostly doing copying that ICE can do (go); if not, stop before training.

\textbf{Main experiment.} GDN hybrids (16 layers, $d{=}768$, ${\approx}140$M parameters, flash-linear-attention), 2.5B tokens (70\% FineWeb-Edu, 30\% code), 4K context, identical data order and schedule across arms:

\begin{center}\small
\begin{tabular}{@{}lll@{}}
\toprule
Arm & Attention & ICE \\
\midrule
A & layer-level 3:1 (4 full-attention layers) & --- \\
B & layer-level 7:1 (2 full-attention layers) & --- \\
C & layer-level 7:1 & input + output, dropout 0.2 \\
D & head-level hybrid at the KV budget of B (HydraHead-style~[7]) & --- \\
E & head-level hybrid at the KV budget of B & input + output, dropout 0.2 \\
F & layer-level 15:1 (1 full-attention layer) & input + output, dropout 0.2 \\
\bottomrule
\end{tabular}
\end{center}

\textbf{Metrics.} (i) Held-out loss split by code/prose and by ``copyable'' tokens (longest match $\ge 4$) vs.\ the rest; (ii) exact-copy tasks: repeat-string, phonebook, MQAR, including 8K--16K extrapolation; (iii) code completion (RepoBench-P, LCC); (iv) needle retrieval with \emph{templated} and \emph{paraphrased} queries, so that gains from pure string matching are not mistaken for retrieval; (v) retrieval-head formation over training checkpoints; (vi) KV memory per token at matched quality. Two seeds for A--C. \textbf{Kill criterion:} C does not beat B on copy tasks beyond seed noise, or ICE hurts paraphrased retrieval and dropout cannot fix it.

\subsection{Compute Requests}

Throughput estimates assume ${\approx}150$--$250$K tokens/s for a 140M model on one H100 (${\approx}2\times10^{18}$ FLOPs, i.e.\ 3--5 hours per run). To leave spare time, each run is budgeted at \textbf{6 H100-hours}, and a 25\% contingency covers reruns, crashed jobs and lower-than-expected throughput. Real throughput will be measured on a short run before any long job.

\begin{center}\small
\begin{tabular}{@{}lrr@{}}
\toprule
Item & H100-hours & Est.\ cost (RunPod) \\
\midrule
Step 0 motivation check (Qwen3.5-0.8B inference) & 6 & \$15--21 \\
Data download \& tokenization (2.5B tokens) & 4 & \$10--14 \\
Arms A--F (6 runs) + second seeds for A--C (3 runs), 6\,h each & 54 & \$135--189 \\
Ablations (input-only / output-only ICE, dropout 0), 6\,h each & 18 & \$45--63 \\
Evaluation & 12 & \$30--42 \\
Setup \& debugging & 10 & \$25--35 \\
Contingency (${\approx}25\%$: reruns, failures, slower throughput) & 26 & \$65--91 \\
\midrule
\textbf{Total} & \textbf{130} & \textbf{\$325--455} \\
\bottomrule
\end{tabular}
\end{center}
Costs assume \$2.5--3.5 per H100-hour. Plus ${\approx}300$\,GB of persistent storage for data and checkpoints. A single GPU suffices; runs are independent, so the work fits in \textbf{2--3 weeks} with time left for reruns, or about one week if 2--3 GPUs run in parallel.

\section{Brief Literature Review}

\textbf{Copying and recall in efficient models.} SSMs fail to copy beyond their state capacity~[1]; associative recall explains most of the attention vs.\ attention-free gap~[4]. Induction heads~[8] and retrieval heads~[5] are the attention circuits that copy; DuoAttention~[6] shows only those heads need full KV.
\textbf{Hybrids.} In GDN/attention hybrids, long-range retrieval lives in the full-attention layers and the recurrent states stay mostly local~[2,3]. HydraHead~[7] keeps only retrieval-critical \emph{heads} as full attention and matches a 3:1 hybrid at 7:1; it still stores KV vectors, whereas ICE adds a vector-free copy path and is tested as a complement (arm E).
\textbf{N-gram mechanisms.} The closest work, Akyürek et al.~[9], hard-wires ``n-gram heads'' (attention over hidden states at matched positions) and lowers perplexity by 2.6--6.7\% at 340M, even for Transformers; it still needs a per-token vector cache and does not aim to remove attention layers. Engram~[10] adds a gated lookup into a \emph{static} n-gram table learned from the corpus; ICE builds its table from the \emph{current context}. Prompt-lookup decoding uses the same matching only to draft tokens at inference. \textbf{Novelty:} to our knowledge, no prior work trains an attention-free, token-ID-only in-context n-gram lookup into a hybrid to lower its full-attention ratio.

\clearpage
\section*{References}
\small
\setlength{\parskip}{4pt}\raggedright

[1] Jelassi et al., \emph{Repeat After Me: Transformers are Better than SSMs at Copying}, arXiv:2402.01032.

[2] \emph{Rethinking the Role of Efficient Attention in Hybrid Architectures}, arXiv:2606.15378.

[3] \emph{Tail-Replay}, arXiv:2608.30310.

[4] Arora et al., \emph{Zoology: Measuring and Improving Recall in Efficient Language Models}, arXiv:2312.04927.

[5] Wu et al., \emph{Retrieval Head Mechanistically Explains Long-Context Factuality}, arXiv:2404.15574.

[6] Xiao et al., \emph{DuoAttention}, arXiv:2410.10819.

[7] Tan et al., \emph{HydraHead: From Head-Level Functional Heterogeneity to Specialized Attention Hybridization}, arXiv:2606.20097.

[8] Olsson et al., \emph{In-context Learning and Induction Heads}, arXiv:2209.11895.

[9] Akyürek et al., \emph{In-Context Language Learning: Architectures and Algorithms}, arXiv:2401.12973.

[10] DeepSeek-AI, \emph{Conditional Memory via Scalable Lookup (Engram)}, arXiv:2601.07372.

\end{document}
